\documentclass[12pt]{article}
\usepackage[T1]{fontenc}
\usepackage[utf8]{inputenc}
\usepackage{lmodern}
\usepackage{textcomp}
\usepackage{enumitem}
\usepackage{geometry}
\usepackage{graphicx}
\usepackage{amsmath, amssymb}
\geometry{margin=1in}
\begin{document}
\begin{center}
Chemistry - Undergraduate Level Assessment 17 \
Total Marks: 40 \
Answer all questions.
\end{center}

\section*{Multiple Choice (10 marks)}
\begin{enumerate}[label=\arabic*.]
\item Which of the following is an n-type semiconductor?
\begin{enumerate}[label=(\alph*)]
    \item Silicon
    \item Diamond
    \item Silicon carbide
    \item Arsenic-doped silicon
\end{enumerate}
\item Calculate the polarization of a proton in a magnetic field of 335 mT and 10.5 T at 298 K.
\begin{enumerate}[label=(\alph*)]
    \item 6.345 x 10^-4 at 0.335 T; 9.871 x 10^-5 at 10.5 T
    \item 0.793 x 10^-4 at 0.335 T; 6.931 x 10^-7 at 10.5 T
    \item 1.148 x 10^-6 at 0.335 T; 3.598 x 10^-5 at 10.5 T
    \item 4.126 x 10^-3 at 0.335 T; 2.142 x 10^-6 at 10.5 T
\end{enumerate}
\item Which of the following is lower for argon than for neon?
\begin{enumerate}[label=(\alph*)]
    \item Melting point
    \item Boiling point
    \item Polarizability
    \item First ionization energy
\end{enumerate}
\item Of the following atoms, which has the lowest electron affinity?
\begin{enumerate}[label=(\alph*)]
    \item F
    \item Si
    \item O
    \item Ca
\end{enumerate}
\item Of the following compounds, which has the lowest melting point?
\begin{enumerate}[label=(\alph*)]
    \item HCl
    \item AgCl
    \item CaCl 2
    \item CCl 4
\end{enumerate}
\end{enumerate}

\section*{True / False (10 marks)}
\textit{Answer True or False}
\begin{enumerate}[label=\arabic*.]
\setcounter{enumi}{5}
\item True or False: In the balanced equation, the ratio of MnO 4- to I- is 6:5.
\item True or False: Raman spectroscopy is a light-scattering technique used to study molecular vibrations.
\item True or False: CH 3 Cl exhibits a 1 H resonance approximately 1.55 kHz downfield from TMS on a 12.0 T NMR spectrometer.
\item True or False: A magnetic field of 6.08 T is optimal for achieving the fastest spin-lattice relaxation in protons of a molecule with a 1 ns correlation time.
\item True or False: The 55 Mn nuclide has three allowed energy levels for its nuclear states.
\end{enumerate}

\section*{Long Form Response (20 marks)}
\textit{Answer each question with a clear and complete explanation.}
\begin{enumerate}[label=\arabic*.]
\setcounter{enumi}{10}
\item Discuss the factors that determine the ratio of line intensities in the EPR spectrum of the t-Bu radical (CH 3)3 C•, and explain the significance of this ratio in understanding radical behavior.
\item Discuss the factors that contribute to the NMR frequency of nuclides, and explain why 17 O exhibits a frequency of 115.5 MHz in a 20.0 T magnetic field.
\end{enumerate}

\vfill
\noindent\textit{End of Paper}
\end{document}